% QUESTAO
Dados o poliedro da figura abaixo

\begin{center}
\vspace{-4mm}
\begin{tikzpicture}[ >=latex, scale=0.5 ]

\coordinate (C) at (0,0);
\coordinate (B) at (-45:-2);
\coordinate (D) at (4,0);
\coordinate (E) at ($(D)+(-45:-2)$);
\coordinate (A) at (80:6);

\draw[thick] (B) -- (C) -- (D);
\draw[dashed] (D) -- (E) -- (B) (A)--(E);
\draw[thick] (A)--(B) (A)--(C) (A)--(D);

\node at ($(A)+(0,0.3)$) {\footnotesize $A$};
\node at ($(B)+(-0.3,0)$) {\footnotesize $B$};
\node at ($(C)+(-135:0.35)$) {\footnotesize $C$};
\node at ($(D)+(-45:0.3)$) {\footnotesize $D$};
\node at ($(E)+(45:0.3)$) {\footnotesize $E$};

\end{tikzpicture}
\vspace{-6mm}
\end{center}

Qual a forma de cada face?

% ALTERNATIVAS
$4$ triangulares e $1$ quadrada
$5$ triangulares
$4$ triangulares
$6$ quadradas
$1$ triangular e $4$ quadradas



